Thiết kế vật lý là mức thiết kế cuối cùng, trong đó mô hình logic được triển khai trên một hệ quản trị cơ sở dữ liệu cụ thể.

Trong đề tài này, hệ quản trị cơ sở dữ liệu được sử dụng là MySQL.

\section{Kiểu dữ liệu và ràng buộc}

Các kiểu dữ liệu được lựa chọn dựa trên đặc điểm của từng thuộc tính.

\begin{table}[htbp]
    \centering
    \caption{Một số kiểu dữ liệu sử dụng}
    \label{tab:data_types}
    \begin{tabularx}{\textwidth}{|p{3cm}|X|p{5cm}|}
        \hline
        \textbf{Kiểu dữ liệu} & \textbf{Mục đích} & \textbf{Ví dụ} \\ \hline
        INT & Khóa định danh, số nguyên & customer\_id \\ \hline
        VARCHAR & Chuỗi ký tự & name, phone, email \\ \hline
        DATETIME & Ngày và thời gian & date\_time \\ \hline
        DATE & Ngày & date \\ \hline
        DECIMAL & Giá trị tiền tệ & payment, amount \\ \hline
        TEXT & Nội dung văn bản dài & diagnosis, treatment \\ \hline
    \end{tabularx}
\end{table}

\begin{table}[htbp]
    \centering
    \caption{Các ràng buộc trong cơ sở dữ liệu}
    \label{tab:constraints}
    \begin{tabularx}{\textwidth}{|p{4cm}|X|}
        \hline
        \textbf{Loại ràng buộc} & \textbf{Mô tả} \\ \hline

        \textbf{PRIMARY KEY} &
        Xác định duy nhất mỗi bản ghi trong một bảng.
        \\ \hline

        \textbf{FOREIGN KEY} &
        Đảm bảo mối liên kết và tính toàn vẹn tham chiếu giữa các bảng.
        \\ \hline

        \textbf{NOT NULL} &
        Yêu cầu thuộc tính phải có giá trị, không được để trống.
        \\ \hline

        \textbf{UNIQUE} &
        Đảm bảo các giá trị trong thuộc tính không bị trùng lặp.
        \\ \hline

        \textbf{DEFAULT} &
        Cung cấp giá trị mặc định cho thuộc tính khi không được chỉ định giá trị.
        \\ \hline

    \end{tabularx}
\end{table}

\section{Dữ liệu mẫu}

Dữ liệu mẫu được xây dựng cho chín bảng nhằm kiểm tra hoạt động của cơ sở dữ liệu và phục vụ các truy vấn ở chương tiếp theo.

Mỗi bảng được khởi tạo một số bản ghi đại diện cho các trường hợp khác nhau, bao gồm booking đã hoàn thành, booking đang chờ xử lý, booking đã xác nhận và booking bị hủy.

\section{Kiểm tra dữ liệu mẫu}

Sau khi dữ liệu mẫu được nhập, có thể kiểm tra số lượng bản ghi bằng truy vấn:

\begin{table}[htbp]
    \centering
    \caption{Các truy vấn đếm số lượng bản ghi}
    \label{tab:count_queries}
    \begin{tabularx}{\textwidth}{|p{7cm}|X|}
        \hline
        \textbf{Truy vấn SQL} & \textbf{Ý nghĩa} \\ \hline

        \texttt{SELECT COUNT(*) AS total\_customers FROM Customer;}
        & Đếm tổng số khách hàng trong bảng Customer.
        \\ \hline

        \texttt{SELECT COUNT(*) AS total\_pets FROM Pet;}
        & Đếm tổng số thú cưng trong bảng Pet.
        \\ \hline

        \texttt{SELECT COUNT(*) AS total\_bookings FROM Booking;}
        & Đếm tổng số lịch đặt khám trong bảng Booking.
        \\ \hline

        \texttt{SELECT COUNT(*) AS total\_invoices FROM Invoice;}
        & Đếm tổng số hóa đơn trong bảng Invoice.
        \\ \hline

        \texttt{SELECT COUNT(*) AS total\_medical\_records FROM MedicalRecord;}
        & Đếm tổng số hồ sơ khám bệnh trong bảng MedicalRecord.
        \\ \hline

    \end{tabularx}
\end{table}

\section{Kết luận chương}

Mô hình logic đã được triển khai ở mức vật lý bằng MySQL với các kiểu dữ liệu và ràng buộc phù hợp. Dữ liệu mẫu được xây dựng để phục vụ việc kiểm tra, truy vấn và đánh giá hệ thống.